\section{Generalization Ability}\label{section6}
In this section, we discuss the propagation of the generalization error between each round of a multi-round sequence generative model.  The generalization ability of an auto-regressive generative language model determines its ability to solve real problems by continuous generation for practical reasoning tasks. We focus on the inter-round propagation of errors when generating a long sequence via \(R\)-rounds, and how the accumulation of these errors acts.

\subsection{The Propagation and Cumulation of Error}

We first consider the case where our model is generated in the \(r\)-th round with an aggregate error due to the dependence on the previously generated contexts. In this process, we consider the effect that the error of the previous round has on the current round.

%\begin{proposition}
%For a function $f$, assume an m step decomposition $f = f_m\circ \cdots \circ f_1: \mathbb R^d \to \mathbb R^d$ , if we define a chain-of-thought $\phi_{i_1}...\phi_{i_k}$ such that a regularization on middle step $\phi_{i_k}$ approximate $f_k$, the required step $m$ for approximation will be reduced from $O(d^2)$ to $O(d\log d)$. 
%\end{proposition}

\begin{lemma}[Aggregate Error]\label{aggregate}
For an auto-regressive generative model over \(R\) rounds, the generalization bound of the aggregate error for each round \(r\) satisfies:
    \[
L_{r}(h_{r}) \leq \sum_{i=1}^{r} \left( \left( \prod_{j=i+1}^{r} \gamma_j \right) \left( \hat{L}_{m,i}(h_{i}) + \epsilon_i \right) \right),
\]
\end{lemma}
where \( L_{r}(h_{r}) = \mathbb{E}_{(x^{(r)}, y^{(r)}) \sim \mathcal{D}^{(r)}} [ \ell(h_{r}(x^{(r)}), y^{(r)}) ] \) is the loss of the model's output sequence in the \(r\)-th round from the expected output. \(\hat{L}_{m,i}(h_{i})\) is the empirical loss computed on \(m\) samples. The \( \gamma_r \) quantifies the impact of errors from round \( r-1 \) on round \( r \). See Appendix~\ref{appendix:aggragate} for proof.
%and \( \lambda_r \) are non-negative weights with \( \sum_{r=1}^R \lambda_r = 1 \).

\begin{theorem}[Cumulative error]\label{cumulative}
For an auto-regressive generative model over \(R\) rounds, the generalization bound of the cumulative error satisfies:
   \[ L(h_R) = \sum_{r=1}^{R}\lambda_r L_{r}(h_{r}) \leq \sum_{i=1}^R \Lambda_i \left( \hat{L}_{m,i}(h_{i}) + \epsilon_i \right).\]
\end{theorem}
where \( \Lambda_i = \sum_{r=i}^R G_{r,i} \) represents the total influence that the generalization error at round \( i \) has on the cumulative error across all subsequent rounds \( r \geq i \), \(\lambda_r\) are non-negative cumulative weights with \(\sum\lambda_r=1\). And here \(
  G_{r,i} = \lambda_r \prod_{j=i+1}^{r} \gamma_j
  \) captures the influence of the generalization error at round \( i \) on the loss at round \( r \). Specifically, \(G_{r,i}\) accounts for the cumulative error for how aggregate errors propagate from round \( i \) through subsequent rounds up to round \( r \). See Appendix~\ref{appendix:cumulative} for proof.
\vspace{-1mm}
\subsection{Cumulative Error Increases as Rounds Grows}

We need to note that, in the real world, models often do not guarantee zero empirical loss (\(\hat{L}_{m,i}(h_{i}) > 0\)) due to limitations of existing optimization methods, model architecture, etc. Let's first assume that the error impact factor \(\gamma = \gamma_r\) is uniform between each round, also a uniform influence factor \(\lambda_r = \lambda, \forall r \in \{1,\cdots, R\}\) and a uniform lower bound \(\eta \geq  \hat{L}_{m,i}(h_{i}) + \epsilon_i \) for simplification. Denote the upper bound of cumulative error \(L(h_R)\) as \( \bar L(h_R) = \sum_{i=1}^R \Lambda_i \left( \hat{L}_{m,i}(h_{i}) + \epsilon_i \right)\). We now observe the trend of the cumulative error upper bound \( \bar L(h_R)\) generated by \(R\)-rounds evolves as the number of rounds \(R\) gradually tends to infinity.

\begin{theorem}[Divergence of Cumulative Error Upper Bound]\label{divergence}
As the generation rounds \(R\) increase to infinity, the upper bound of generation cumulative error \(L(h_R)\) satisfies:
    \[
\lim_{R \to \infty} \bar L(h_R) =\lim_{R \to \infty} \frac{\eta \lambda}{1 - \gamma} \sum_{i=1}^R \left( 1 - \gamma^{R - i + 1} \right) \to \infty.
\]
\end{theorem}

Theorem~\ref{divergence} implies that there is no supremum on the cumulative error as rounds \(R\) increases. This means that, in reality, considering models with limited accuracy and limited training, we do not always get the results we expect using multiple rounds of inference. This may sound like a disappointing result, suggesting that if the model is allowed to keep generating round by round, its final generated content can be uncontrollable. But the good news is that, for most practical scenarios, the rounds are typically finite, hence the cumulative error is within an acceptable range. This suggests that even though the content the auto-regressive language model generates in a finite number of rounds may not be the precise solution to the problem that we expect, it is at least to a degree relevant to a specific topic that we expect. See Appendix~\ref{appendix:divergence} for proof of Theorem~\ref{divergence}.
 %\jx{Q: the previous sentence is a bit confusing. check}.

\subsection{Multi-Round Generation Techniques as Intervention}

We now consider the use of techniques, such as Chain-of-Thought~\citep{wei2022chain} or self-correction~\citep{madaan2023selfrefine}, and multi-agent dialog~\citep{pmlr-v235-khan24a} to intervene in the generation process of language models. Chain-of-thought restricts the intermediate steps in multiple rounds of generation by means of a prompt template of the solution idea, in order to block the effect of errors in this round on the subsequent generation. Self-correction limits the effect of errors by constantly making multiple small adjustments to the results obtained in the previous round. Multi-agent dialog, on the other hand, limits the aggregate error by the information given by other agents, thus reducing the error propagation. It is worth noting that these tricks are ultimately a kind of generative process intervention to reduce \(\gamma_r\) to \(\gamma'\) at certain rounds.

\begin{theorem}\label{intervention} If we allow intervention several times during the generation process by making \(\gamma_r\) decrease to \(\gamma'\), then let \( h_{i,r}\) be the number of hint rounds between \( i+1 \) and \( r \). It can be shown that
    the reduction in cumulative error can be given by:
  \[
  \Delta L(h_R) = \sum_{i=1}^{R} \left( 1 - \kappa_i \right) \Lambda_i \left( \hat{L}_{m,i}(h_i) + \epsilon_i \right)
  \]
 where \(\kappa_i =  \mathbb{E}_{r \sim \mu_i} \left[ \left( \frac{\gamma'}{\gamma} \right)^{h_{i,r}} \right] \), \( \mu_i \) is a probability distribution over \( r \) defined by: \(\mu_i(r) = \frac{\lambda_r \gamma^{r - i}}{\Lambda_i}\).
\end{theorem}

Theorem~\ref{intervention} means that if we obtain a good intervention such that the error propagation at certain rounds is effectively controlled by a smaller amount to \(\gamma'\), the cumulative error of sequence generation will be effectively controlled. To achieve effective control over the cumulative error, i.e. large cumulative error reduction \(\Delta L(h_R)\), we expect a smaller \(\kappa_i\). This can be accomplished in two ways: (1) Improving the quality of hints: A good hint can effectively block error propagation in a single-round generation, leading to a smaller \(\gamma'\), eventually smaller \(\kappa_i\). (2) Increase the number of hints: By increasing the number of hints between round \(i\) to \(r\), we will get a larger \(h_{i,r}\), which ultimately leads to a decrease in \(\kappa_i\) in average regarding \(i\). See Appendix~\ref{appendix:intervention} for proof of Theorem~\ref{intervention}.


